\section{System Overview: The End-to-End Pipeline}
\label{sec:overview}

Before we dissect any single stage, it helps to see the whole machine at once. Cadence
takes one stream of numbers --- the wobble of a single accelerometer strapped to the
ankle --- and turns it, every two seconds, into one of three plain-English words:
\textcolor{grn}{\textbf{STILL}}, \textcolor{steel}{\textbf{WALKING}}, or
\textcolor{freeze}{\textbf{FREEZE}}. That verdict is shown on the body (a small OLED and
a ring of coloured LEDs) and, optionally, streamed over Wi-Fi to a live laptop dashboard.
Everything else in this report is an expansion of one of the boxes in
Figure~\ref{fig:pipeline}.

\begin{keyidea}
Cadence is a \textbf{monitor}, not a therapy. It \emph{detects} and \emph{displays}
freeze-of-gait episodes so they can be logged and reviewed; it never tries to cue, buzz,
or treat the wearer. A coin vibration motor is physically present on the build but is
\textbf{disabled in firmware} --- the only outputs are visual and on-screen.
\end{keyidea}

\subsection{The journey of one window}

The pipeline is a straight line with a single branch. Signal flows left to right; the
two analysis paths in the middle answer the same question --- \emph{is this a freeze?}
--- in two very different ways.

\begin{figure}[H]
\centering
\resizebox{0.985\textwidth}{!}{%
\begin{tikzpicture}[
    font=\sffamily,
    >={Stealth[length=2.6mm]},
    block/.style={rounded corners=2pt,draw,line width=0.9pt,minimum height=12mm,
                  minimum width=20mm,align=center,fill=white,inner sep=5pt},
    sense/.style={block,draw=teal,fill=paleteal,text=ink},
    compute/.style={block,draw=navy,fill=palenavy,text=ink},
    decide/.style={block,draw=freeze,fill=palefreeze,text=ink},
    output/.style={block,draw=grn,fill=white,text=ink},
    flow/.style={->,line width=1.0pt,draw=mute},
    lbl/.style={font=\scriptsize\sffamily,text=mute,align=center},
  ]
  % local sub-label size: one step up from \footnotesize for print legibility
  \newcommand{\sub}{\small}

  % --- main spine ---------------------------------------------------------
  \node[sense]   (acc)  at (0,0)
        {\textbf{Accelerometer}\\[1pt]{\sub LIS3DH, 3 axis}\\{\sub 64\,Hz, $\pm$8\,g}};
  \node[sense]   (mag)  at (3.5,0)
        {\textbf{Magnitude}\\[1pt]{\sub $\norm{a}$, mean-}\\{\sub removed (DC out)}};
  \node[sense]   (win)  at (7.0,0)
        {\textbf{Sliding window}\\[1pt]{\sub 256 samp (4\,s)}\\{\sub 50\% overlap}};

  % branch hub
  \node[compute] (fi)   at (11.0,1.7)
        {\textbf{Welch PSD}\\{\sub band power}\\{\sub $\boldsymbol{\to}$ \textbf{Freeze Index}}};
  \node[compute] (cnn)  at (11.0,-1.7)
        {\textbf{FoGNet (CNN)}\\{\sub 1-D conv, $\sim$18k par}\\{\sub $\to P(\text{freeze})$}};

  \node[decide]  (gate) at (15.0,0)
        {\textbf{Energy gate}\\[1pt]{\sub movement}\\{\sub energy $>$ floor?}};
  \node[decide]  (deb)  at (18.5,0)
        {\textbf{Debounce}\\[1pt]{\sub assert / release}\\{\sub after 2 windows}};
  \node[decide]  (state) at (22.0,0)
        {\textbf{State}\\[1pt]{\sub STILL / WALK}\\{\sub / FREEZE}};

  % outputs
  \node[output]     (oled) at (25.6,1.4)
        {\textbf{On body}\\{\sub OLED + NeoPixels}};
  \node[output]     (dash) at (25.6,-1.4)
        {\textbf{Dashboard}\\{\sub Wi-Fi $\to$ laptop}};

  % --- arrows -------------------------------------------------------------
  \draw[flow] (acc)  -- (mag);
  \draw[flow] (mag)  -- (win);
  % win -> branch
  \draw[flow] (win.east) -- ++(0.9,0) |- (fi.west);
  \draw[flow] (win.east) -- ++(0.9,0) |- (cnn.west);
  % branch -> gate  (FI is the shipped path: solid; CNN: dashed)
  \draw[flow] (fi.east)  -| ($(gate.north)+(0,0.0)$);
  \draw[flow,dashed,draw=navy] (cnn.east) -| ($(gate.south)+(0,0.0)$);
  \draw[flow] (gate) -- (deb);
  \draw[flow] (deb)  -- (state);
  \draw[flow] (state.east) -- ++(0.8,0) |- (oled.west);
  \draw[flow] (state.east) -- ++(0.8,0) |- (dash.west);

  % --- group brackets / stage labels -------------------------------------
  \begin{scope}[on background layer]
    \node[fit=(acc)(mag)(win),draw=teal,rounded corners,line width=0.6pt,
          dashed,inner sep=6pt] (gsense) {};
    \node[fit=(fi)(cnn),draw=navy,rounded corners,line width=0.6pt,
          dashed,inner sep=6pt] (gcomp) {};
    \node[fit=(gate)(deb)(state),draw=freeze,rounded corners,line width=0.6pt,
          dashed,inner sep=6pt] (gdec) {};
  \end{scope}
  \node[lbl,text=teal,above=2pt of gsense] {\textbf{1.\ SENSE} \,(teal)};
  \node[lbl,text=navy,above=2pt of gcomp]  {\textbf{2.\ ANALYSE} \,(navy)};
  \node[lbl,text=freeze,above=2pt of gdec] {\textbf{3.\ DECIDE} \,(red)};

  % shipped vs upgrade annotation
  % placed just outside the dashed ANALYSE box, clear of its border, with a
  % white fill so the labels never collide with the dashed boundary.
  \node[lbl,text=navy,fill=white,inner sep=1.5pt,
        anchor=south west] at ($(fi.north east)+(0.18,0.30)$)
        {\textbf{shipped} (glass box)};
  \node[lbl,text=navy,fill=white,inner sep=1.5pt,
        anchor=north west] at ($(cnn.south east)+(0.18,-0.30)$)
        {\emph{upgrade} (not shipped)};

\end{tikzpicture}}
\caption{The Cadence pipeline, end to end. One accelerometer stream is reduced to a
\textsc{still}/\textsc{walking}/\textsc{freeze} verdict every 2\,s. Stages are colour
coded: \textcolor{teal}{\textbf{sensing}} (teal), \textcolor{navy}{\textbf{analysis}}
(navy), \textcolor{freeze}{\textbf{decision}} (red). The \emph{shipped} detector is the
glass-box Freeze Index (solid path); the 1-D CNN \textsc{FoGNet} is a drop-in
\emph{upgrade} (dashed) explored later, not on the deployed board.}
\label{fig:pipeline}
\end{figure}

Reading Figure~\ref{fig:pipeline} from left to right:

\begin{enumerate}
  \item \textbf{Sense.} The onboard LIS3DH MEMS accelerometer samples three axes at
    \SI{64}{\hertz} over $\pm\SI{8}{\gram}$. The wide range matters: a heel-strike
    momentarily exceeds $\pm\SI{2}{\gram}$, so a $\pm\SI{2}{\gram}$ part would clip.
    The Nyquist frequency of \SI{32}{\hertz} comfortably covers every band we care about
    (the highest, freeze, tops out at \SI{8}{\hertz}).
  \item \textbf{Magnitude.} The three axes collapse into one orientation-invariant
    scalar $\norm{a}=\sqrt{a_x^2+a_y^2+a_z^2}$; subtracting its mean removes gravity and
    any DC offset, leaving pure motion to analyse.
  \item \textbf{Window.} The stream is chopped into overlapping
    \SI{256}{\text{-sample}} (\SI{4}{\second}) windows with a \SI{128}{\text{-sample}}
    (\SI{2}{\second}) hop. Fifty-percent overlap yields one fresh decision every
    \SI{2}{\second}.
  \item \textbf{Analyse (the branch).} Each window is scored by either the shipped
    \emph{Freeze Index} --- a ratio of spectral band powers --- or, as an upgrade, the
    \textsc{FoGNet} CNN that reads the raw window directly. Both answer the same yes/no
    question.
  \item \textbf{Decide.} A movement-energy \emph{gate} forces \textsc{still} when the leg
    is barely moving (so noise cannot masquerade as a freeze); a two-window
    \emph{debounce} stops a single bump from flipping the verdict; the result is the
    committed state.
  \item \textbf{Display.} The state lights the body-worn OLED and NeoPixel ring and, on
    the optional networked build, drives the live laptop dashboard.
\end{enumerate}

\begin{intuition}
The whole device is one long funnel. A 3-axis, 64-per-second firehose of numbers is
squeezed, stage by stage, down to a single word twice a second. Each stage throws away
information that does not help answer ``are the feet trembling in place?'' --- first
\emph{direction} (magnitude), then \emph{time detail} (the spectrum), then \emph{ambiguity}
(the gate and debounce).
\end{intuition}

\subsection{Two detectors, one question}

The single most important thing to understand about Cadence is that the box labelled
\textbf{Analyse} contains two interchangeable answers to the same question, and the report
treats them very differently.

\begin{itemize}
  \item The \textbf{Freeze Index (FI)} is the \emph{shipped} detector --- a transparent
    ``glass box''. It is the ratio of energy in the freeze band (\SIrange{3}{8}{\hertz})
    to energy in the locomotor band (\SIrange{0.5}{3}{\hertz}). When the feet shuffle in
    place, the freeze band lights up and the ratio jumps. Every step of this calculation
    is inspectable, which is exactly why it runs on the board.
  \item \textsc{FoGNet}, a small 1-D \textbf{convolutional neural network}
    ($\sim$18k parameters), is the ``given more time'' \emph{upgrade}. It learns its own
    features from the raw $3\times256$ window. It is \textbf{not} on the deployed device;
    it appears later as a benchmark of what a learned model buys (and costs) over the
    glass box.
\end{itemize}

\begin{takeaway}
The FI path is what actually ships and runs on the ankle; the CNN is an offline
comparison. Wherever this report says ``the detector'' without qualification, it means the
glass-box Freeze Index.
\end{takeaway}

Figure~\ref{fig:realsignal} shows the pipeline's raw input and its headline statistic on a
real \SI{41.4}{\second} capture. The signal is calm during standing and walking, erupts
into rapid trembling during the freeze phase, and the Freeze Index --- computed exactly as
the device computes it --- tracks that eruption faithfully.

\begin{figure}[H]
\centering
\includegraphics[width=0.86\linewidth]{capture1_trace}
\caption{Real input and output of the analysis stage on \code{capture1.csv}
(\SI{64}{\hertz}, \SI{41.4}{\second}). \emph{Top:} the mean-removed accelerometer
magnitude through four labelled phases. The freeze phase (centre) shows the tell-tale
rapid \SIrange{3}{8}{\hertz} trembling in place. \emph{Bottom:} the Freeze Index over the
same window; it stays near 1 during stillness and walking and spikes above 20 during the
freeze, far past the \num{1.815} trip line.}
\label{fig:realsignal}
\end{figure}

\subsection{Three boards, three jobs}

Cadence runs untethered, but the optional networked build spreads the work across three
physical boards (Table~\ref{tab:boards}). The simplest deployment is just the first board:
the Circuit Playground Express senses, runs the on-board FI detector, and drives its own
OLED and LEDs --- no laptop required. Adding the ESP32 and a laptop only changes \emph{where}
the verdict is displayed, never how it is computed.

\begin{table}[H]
\centering
\small
\renewcommand{\arraystretch}{1.25}
\begin{tabular}{@{}>{\bfseries}l l l@{}}
\toprule
Board & Role & Job in the pipeline \\
\midrule
CPX & Sensor + detector & Samples the LIS3DH at \SI{64}{\hertz}, computes magnitude, the \\
    & (always present)  & Freeze Index, the gate and debounce; drives the OLED + NeoPixels. \\
ESP32 & Wi-Fi bridge & Relays the stream over \mbox{UART\,$\to$\,Wi-Fi/TCP} to the laptop. \\
      & (optional)   & Carries data only; makes no decision. \\
Laptop & Dashboard host & Runs the asyncio + WebSocket server that renders the live \\
       & (optional)     & display panels and, offline, trains/evaluates \textsc{FoGNet}. \\
\bottomrule
\end{tabular}
\caption{The three boards. Only the \textbf{CPX} is required; the freeze decision always
happens there. The ESP32 and laptop add a live remote dashboard but do not move the
detector off the ankle.}
\label{tab:boards}
\end{table}

A matching split exists in firmware. One role, the \emph{recorder}
(\code{cpx\_fog\_logger.ino}), logs auto-labelled accelerometer data to the CPX's
\SI{2}{\mega\byte} flash for later analysis; the other, the \emph{monitor}
(\code{cpx\_fog\_standalone.ino}), runs the live OLED-and-LED freeze display. The figures
of real data throughout this report --- like Figure~\ref{fig:realsignal} --- come from the
recorder; the device the reader would wear runs the monitor.

\begin{keyidea}
The decision lives on the ankle. Whether or not a laptop is in the loop, the CPX alone
turns raw acceleration into a \textsc{freeze} verdict and shows it on the body. The rest of
the report simply opens up each box in Figure~\ref{fig:pipeline}: the spectral maths
(\S\ref{sec:overview}\,$\to$ Fourier), the Freeze Index and its threshold, the
decision logic, and the CNN upgrade.
\end{keyidea}
